\section{Research Problem and Question}
% Define the gap and state the research question and sub-questions.

\citet{hsu2026rashomon} measure the leakage of a Rashomon release as the average error of one reconstruction against uniform guessing. A single global error hides what a disclosure policy needs to know: how much depends on what the attacker already knows, and which records are exposed. Our question is: \emph{when $m$ sparse decision trees from one Rashomon set are released with their leaf class counts, how well can an attacker reconstruct the sex of the training records, and can the attacker tell which reconstructed records are correct?}
\begin{description}[leftmargin=1.2em,itemsep=0pt,topsep=2pt]
  \item[SQ1.] How much does knowing marital status, education or age for every row lower the error, compared with releasing more trees, and does it matter which attribute is known?
  \item[SQ2.] Does it matter whether the released trees are a random subset of the Rashomon set or the first $m$ in TreeFARMS's enumeration order?
  \item[SQ3.] Can the attacker give each reconstructed record a confidence that is calibrated without fitting on ground truth, using only the release and a public sample, and how much public data does that need?
\end{description}
SQ2 came up during the project, when the first trees in enumeration order turned out to be near-copies (Section~\ref{sec:selection}). We do not propose a new attack: we adapt DRAFT and add repeated solving guided by public data, close in spirit to \citet{dick2023confidence}.
